\chapter{Энтропийно-ориентационный переход вещественного семейного состояния}

Предыдущая проверка свела задачу рождения $\mathbb{RP}^2$ к расщеплению
изотропного спектра $3\to1+2$. Теперь проверяется, способен ли на это уже
имеющийся вариационный принцип состояния и какой минимальный нелинейный
член действительно создаёт такую фазу.

\section{Линейный функционал Гиббса не нарушает симметрию}

Все завершённые проектные вариации состояния имеют базовую форму
\begin{equation}
 \mathcal F_{H,\beta}(R)
 =\Tr(RH)+\beta^{-1}\Tr(R\log R),
 \qquad R\in\mathcal D_{\mathbb R,3}.
 \label{eq:v6-linear-gibbs-family}
\end{equation}
При конечном $\beta>0$ её единственный минимум равен
\begin{equation}
 R_\beta=\frac{e^{-\beta H}}{\Tr e^{-\beta H}}.
 \label{eq:v6-linear-gibbs-minimum}
\end{equation}

Если $H$ инвариантен относительно неприводимого действия $SO(3)$ на
$\mathbb R^3$, то по лемме Шура $H=cI_3$ и
$R_\beta=I_3/3$. Если $H$ не скалярен, он заранее фиксирует собственные
оси, и минимум не имеет свободной орбиты $\mathbb{RP}^2$. В пределе
$\beta\to\infty$ можно получить проектор основного состояния, но его ось
опять задана матрицей $H$, а конечного фазового перехода нет.

Проектный $R_4^+$ имеет три различных собственных значения
\begin{equation}
 12.608881\ldots,quad14.758684\ldots,quad20.632435\ldots.
\end{equation}
Он беспараметрически выбирает одну фиксированную семейную ось и полезен для
ориентации ветви, но не создаёт вырожденное вакуумное многообразие.

\begin{theorem}[Линейный Gibbs запрет]
Конечномерный функционал \eqref{eq:v6-linear-gibbs-family} не может
одновременно сохранить $SO(3)$-симметрию родителя и иметь множество
минимумов $\mathbb{RP}^2$. Симметричный $H$ оставляет изотропный минимум,
а несимметричный $H$ выбирает ось явно.
\end{theorem}

\section{Минимальный нелинейный контроль}

Простейший вращательно-инвариантный функционал, конкурирующий с энтропией,
имеет вид
\begin{equation}
 \mathcal F_\kappa(R)
 =\Tr(R\log R)
 -\kappa\left(\Tr R^2-\frac13\right),
 \qquad \kappa\ge0.
 \label{eq:v6-entropy-alignment-functional}
\end{equation}
Это конечномерный аналог энтропийно-ориентационной конкуренции
Майера--Заупе. Для одноосного семейства \eqref{eq:v6-uniaxial-mixed-state}
получаем
\begin{align}
 f_\kappa(x)
 &=x\log x+(1-x)\log\frac{1-x}{2}\notag\\
 &\quad-\kappa\left(x^2+\frac{(1-x)^2}{2}-\frac13\right).
 \label{eq:v6-uniaxial-free-energy}
\end{align}
Условие стационарности равно
\begin{equation}
 \log\frac{2x}{1-x}=\kappa(3x-1).
 \label{eq:v6-uniaxial-stationarity}
\end{equation}

Это сечение не является догадкой о форме минимума. Функционал зависит
только от собственных значений $p_i$ состояния. Во внутренней стационарной
точке
\begin{equation}
 \log p_i-2\kappa p_i=\text{const}.
\end{equation}
Функция слева имеет единственную точку поворота, поэтому среди трёх
$p_i$ возможно не более двух различных значений. Граница симплекса при
конечном $\kappa$ не является минимумом из-за сингулярного энтропийного
наклона. Следовательно, всякая неизотропная глобальная стационарная точка
имеет одноосный спектр, и проверка \eqref{eq:v6-uniaxial-free-energy}
охватывает полный симплекс состояний.

При
\begin{equation}
 \boxed{\kappa_c=\log4}
 \label{eq:v6-exact-ordering-threshold}
\end{equation}
изотропный минимум $x=1/3$ имеет ту же энергию, что и одноосный минимум
$x=2/3$. Оба локально устойчивы, поскольку изотропная точка теряет
локальную устойчивость только при $\kappa=3/2>\log4$. Следовательно,
переход первого рода создаёт орбиту
\begin{equation}
 R=\frac23P+\frac16(I_3-P),\qquad P\in\mathbb{RP}^2,
 \label{eq:v6-critical-rp2-state}
\end{equation}
ещё до полной очистки.

Этот результат согласуется с признанной физикой изотропно-нематического
перехода Майера--Заупе и с энтропийным потенциалом Ball--Majumdar; см.
\path{arXiv:0808.1870}, Ball--Majumdar, Mol. Cryst. Liq. Cryst. 525
(2010), и \path{arXiv:2002.11855}.

\begin{proposition}[Точный минимальный фазовый механизм]
Отрицательный квадратичный инвариант чистоты вместе с энтропией достаточен
для конечного перехода $I_3/3\to\mathbb{RP}^2$. Критический безразмерный
коэффициент в нормировке \eqref{eq:v6-entropy-alignment-functional} равен
$\log4$.
\end{proposition}

\section{Даёт ли проект требуемый знак и коэффициент}

Нет. Вариационный проверка семейного состояния содержит линейный член $\Tr(RR_4)$,
а не $-\Tr R^2$. Ретроспективная программа TOE~6.5 оставила
$\Gamma_{\rm fluc}[R,M]$ невычисленным. Корреляционная чистота
$\Tr R^2$ использовалась как диагностическая величина, но знак её вклада
в действие не был выведен.

Топологический отклик $1/7$ также не закрывает коэффициент. Он вычислен
после выбора дефектного проектора, положителен и зависит от отдельной
щели. Его нельзя отождествлять с $\kappa$ без операторного отображения и
общей нормировки. Даже формальная подстановка $\kappa=1/7$ оставила бы
изотропную фазу, поскольку $1/7<\log4$.

Наконец, каноническое действие обменного моста
$2(1-\lambda^2)^2/7$ стабилизирует замкнутый мост. Поэтому ни один
готовый член текущего родителя не даёт требуемого отрицательного
ориентационного взаимодействия.

\begin{nogocriterion}[Запрет назначения $\kappa$]
Нельзя объявлять \eqref{eq:v6-entropy-alignment-functional} родительским
действием, пока $-\kappa\Tr R^2$, его знак и нормировка не получены
интегрированием уже существующей степени свободы или точным следовым
тождеством. Числа $1/7$, $15$ и $\log4$ нельзя отождествлять по
численному сходству или перенормировке задним числом.
\end{nogocriterion}

\section{Что найдено и что проверять дальше}

Найден точный рабочий механизм, но ещё не его проектный источник. Он уже
не требует внешней оси, чистого проектора или добавления новой размерности.
Нужен один отрицательный коллективный инвариант состояния. Наиболее
экономный следующий тест: интегрировать обменный Real-мост или его
флуктуации при общем семейном состоянии $R$ и проверить, возникает ли
$-\Tr R^2$ с фиксированным коэффициентом.

\begin{protocol}
\begin{enumerate}
 \item линейный Gibbs functional при конечной температуре:
 \textbf{единственный полный Gibbs-минимум};
 \item симметричный линейный Hamiltonian:
 \textbf{оставляет $I_3/3$};
 \item несимметричный линейный Hamiltonian:
 \textbf{вставляет ось явно};
 \item проектный $R_4^+$:
 \textbf{выбирает фиксированную ось, не $\mathbb{RP}^2$};
 \item минимальный нелинейный механизм:
 \textbf{энтропия минус $\kappa\Tr R^2$};
 \item точный порог первого рода:
 \textbf{$\kappa_c=\log4$};
 \item состояние при сосуществовании:
 \textbf{спектр $(2/3,1/6,1/6)$};
 \item вакуумная орбита:
 \textbf{$\mathbb{RP}^2$};
 \item знак и коэффициент $\kappa$ из текущего родителя:
 \textbf{не выведены};
 \item рождение пространства и материи:
 \textbf{ещё не доказано};
 \item следующая проверка:
 \textbf{индуцирует ли обменный мост ориентационное взаимодействие}.
\end{enumerate}
\end{protocol}

Машинный сертификат сохранён в
\path{s2t_v6_real_qutrit_purification_transition_gate_results.json}.